\begin{textsample}{POS Dim 3 – rr – Score 19.00 – rr/rr\_inf\_17.txt}  \label{ex:f3_pos_224}
3.6.4 AS ATIVIDADES PROFISSIONAIS NA \textbf{INSTITUIÇÃO} DE EDUCAÇÃO \textbf{INFANTIL} A \textbf{instituição} de Educação \textbf{Infantil} deve ser considerada como um ambiente familiar, afetuoso contendo suas normas de funcionamento conforme o entendimento de sua comunidade para melhor atender as \textbf{crianças} oferecendo o bem-estar propiciando um lugar seguro com ações que contribuam para a qualidade do atendimento oportunizando -as experiências, aprendizagens e que \textbf{brinquem} e convivam com a diversidade humana, cultural e social.
Para que se tenham as diversidades fortalecidas, enriquecidas, faz -se necessário que as famílias acompanhem as vivências e as produções das \textbf{crianças}, pois, a parceria estabelecida entre escola e família é primordial para um trabalho eficaz nas \textbf{instituições} de Educação \textbf{Infantil}.
Essa participação da família na escola permite a proximidade com os profissionais e, consequentemente, troca de conhecimentos promovendo melhores elementos para \textbf{apoiar} as \textbf{crianças} nas suas vivências e assim se apropriarão ainda mais de suas potencialidades, seus gostos e dificuldades, isto, contribuirá no processo do \textbf{cuidar} e educar.
As \textbf{crianças} possuem diferenças de temperamento, atitudes, credo, gênero, etnia, características físicas, habilidades e de conhecimentos, por isso, deve -se criar situações de aprendizagem em que a questão da diversidade seja abordada nessas \textbf{instituições} ( BRASIL, 1998 ).
De acordo com as Diretrizes Curriculares Nacionais, Resolução de no 09/2009, as práticas pedagógicas da Educação \textbf{Infantil} por meio dos eixos estruturantes interações e \textbf{brincadeiras} devem garantir as experiências que : - Promovam o conhecimento de si e do mundo por meio da ampliação de experiências \textbf{sensoriais}, expressivas, corporais que possibilitem \textbf{movimentação} ampla, expressão da individualidade e respeito pelos ritmos e \textbf{desejos} da \textbf{criança} ; - Favoreçam a \textbf{imersão} das \textbf{crianças} nas diferentes linguagens e o progressivo domínio por elas de vários gêneros e formas de expressão : gestual, verbal, \textbf{plástica}, \textbf{dramática} e musical ; - Possibilitem às \textbf{crianças} experiências de narrativas, de apreciação e interação com a linguagem oral e escrita, e convívio com diferentes suportes e gêneros textuais orais e escritos ; - Recriem, em contextos significativos para as \textbf{crianças}, relações quantitativas, medidas, formas e orientações espaço temporais ; - Ampliem a confiança e a participação das \textbf{crianças} nas atividades individuais e coletivas ; - Possibilitem situações de aprendizagem mediadas para a elaboração da autonomia das \textbf{crianças} nas ações de \textbf{cuidado} pessoal, auto-organização, saúde e bem-estar ; - Possibilitem vivências éticas e estéticas com outras \textbf{crianças} e grupos culturais, que alarguem seus padrões de referência e de identidades no diálogo e conhecimento da diversidade ; - Incentivem a \textbf{curiosidade}, a exploração, o encantamento, o questionamento, a indagação e o conhecimento das \textbf{crianças} em relação ao mundo físico e social, ao tempo e à natureza ; - Promovam o relacionamento e a interação das \textbf{crianças} com diversificadas manifestações de música, artes \textbf{plásticas} e gráficas, cinema, fotografia, dança, teatro, poesia e literatura ; - Promovam a interação, o \textbf{cuidado}, a preservação e o conhecimento da biodiversidade e da sustentabilidade da vida na Terra, assim como o não desperdício dos recursos naturais ; - Propiciem a interação e o conhecimento pelas \textbf{crianças} das manifestações e tradições culturais brasileiras ; - Possibilitem a utilização de gravadores, projetores, computadores, máquinas \textbf{fotográficas}, e outros recursos tecnológicos e midiáticos.
O profissional que está na Educação \textbf{Infantil} tem a responsabilidade de garantir os direitos de aprendizagem das \textbf{crianças} assegurados nos eixos estruturantes, pois, deve proporcionar as experiências que contribuam no desenvolvimento das capacidades cognitivas, como atenção, memória raciocínio e o bem-estar em ambiente \textbf{cheio} de pluralidade.
Para isso, a organização da prática pedagógica nas unidades de Educação \textbf{Infantil} deve favorecer o interesse da \textbf{criança}.
É importante destacar que, > [... ] práticas pedagógicas devem ocorrer de modo a não fragmentar a \textbf{criança} nas suas possibilidades de viver experiências, na sua compreensão do mundo feita pela totalidade de seus sentidos, no conhecimento que constrói na relação intrínseca entre a razão e emoção, expressão corporal e verbal, experimentação prática e elaboração conceitual.
As práticas envolvidas nos atos de alimentar -se, tomar banho, trocar fraldas e controlar os esfíncteres, na escolha do que vestir, na atenção aos riscos de adoecimento mais fácil nessa faixa \textbf{etária}, no âmbito da Educação \textbf{Infantil} não são apenas práticas que respeitam o direito da \textbf{criança} de ser bem atendida nesses aspectos, como cumprimento do respeito a sua dignidade como pessoa humana.
Elas são também práticas que respeitam e atendem ao direito da \textbf{criança} de apropriar -se por meio de experiências corporais, dos modos estabelecidos culturalmente de alimentação e promoção de saúde, de relação com o próprio corpo e \textbf{consigo} mesma, mediada pelas professoras e professores que intencionalmente planejam e \textbf{cuidam} da organização dessas práticas ( DCNEI‘s 2009, p.9 - 10 ).
Nesse prisma, toda ação do trabalho pedagógico nas \textbf{instituições} de ensino de educação \textbf{infantil} deve ser pensada no interesse da \textbf{criança}, propiciando a integração dos aspectos físicos, emocionais, afetivos, cognitivo/linguísticos e sociais, integrando o \textbf{cuidar}, o educar e o \textbf{brincar}.

% matched lemmas: apoiar, brincadeira, brincar, cheio, conseguir, criança, cuidado, cuidar, curiosidade, desejo, dramático, etário, fotográfico, imersão, infantil, instituição, movimentação, plástico, sensorial
\end{textsample}
